\documentclass[11pt, oneside]{article}   	% use "amsart" instead of "article" for AMSLaTeX format
\usepackage{geometry}                		% See geometry.pdf to learn the layout options. There are lots.
\geometry{letterpaper}                   		% ... or a4paper or a5paper or ... 
%\geometry{landscape}                		% Activate for rotated page geometry
%\usepackage[parfill]{parskip}    		% Activate to begin paragraphs with an empty line\usepackage{graphicx}				% Use pdf, png, jpg, or eps§ with pdflatex; use eps in DVI mode
								% TeX will automatically convert eps --> pdf in pdflatex
								
%\usepackage{{Gratzer-Color-Scheme}	% Active to color theorems red and lemmas blue
%\raggedright	
\usepackage{setspace}
\setstretch{1.1}

\usepackage{amssymb}

%SetFonts

%SetFonts


\title{The Pennsyvalnia State University \\
 \Large \& The Briseño Lab Introductory Manual}
\author{Ricardo I. Alcalá Briseño}
\date{}							% Activate to display a given date or no date

\begin{document}
\maketitle

\section{Introduction}
Welcome to The Pennsylvania State University (PSU) and the Department of Plant Pathology and Environmental Microbiology (PPEM). Penn State has multiple campuses across Pennsylvania; PPEM is located at the University Park campus in State College, Centre County, Pennsylvania.

Your office address is:
Your Office #, Buckhout Laboratory, University Park, PA 16802.
The official 911 address is:
361 Science Drive, University Park, PA 16802.

I will provide several helpful pointers to support your arrival and initial transition. I may miss some steps; please feel free to add notes or suggestions, and I will incorporate them into this manual.

The PPEM Department is housed in Buckhout Laboratory. A helpful resource for navigating campus is the Penn State Maps website (https://www.psu.edu/maps), which provides locations for buildings, offices, and classrooms. Buckhout Laboratory is located near the HUB-Robeson Center (“the HUB”), which offers dining options, coffee, study spaces, printing services, and general student resources.

My office is located on the first floor of Buckhout Laboratory (Room 101). Our primary lab space is Room 107, and we are currently sharing temporary space in Lab 106. Your office will be on the third floor in the graduate student office area; your assigned desk will be either Desk 19 or 21. Please choose thoughtfully (e.g., proximity to windows or lighting), as this will be your workspace for the next several years.

You will need to coordinate with Ed Kaiser (eqk3@psu.edu) to obtain keys for offices and laboratories. Ed is responsible for the spawn laboratory, located on the first floor (Room 117). You should also contact Kim Hall (khall@psu.edu), the Graduate Student Assistant (Office 210), who will have access to your computer and can assist with questions about classes, credits, and graduate program logistics. Linda Nóbrega (lkn5169@psu.edu), the Administrative Support Assistant (Office 211), can help with general administrative and departmental matters.


\section{Beginning at PSU}

\subsection{Internet}
Upon arrival, you can temporarily connect to \texttt{PSU Guest} WiFi. This network provides limited access and should be used only to complete the initial setup. For full internet access, Penn State uses the \texttt{eduroam} network, the primary and secure WiFi service across campus.

To set up eduroam, visit https://www.it.psu.edu/get-connected/ and select \textit{Connect to WiFi}. Follow the instructions for your device (laptop, phone, or tablet). You will be prompted to install Penn State–provided configuration software or certificates; this is required for secure access.

Access to eduroam and most Penn State services requires two-factor authentication (2FA). PSU uses the Microsoft Authenticator app. Install the app on your smartphone using the instructions at https://mobile.psu.edu/project/microsoft-authenticator/. Once 2FA is enabled, you will be able to log in to Canvas, LionPATH, email, and other PSU systems.

If you experience issues connecting to WiFi or setting up authentication, contact Penn State IT support or visit an on-campus IT service desk.

\subsection{Computer set up}
Penn State–issued computers must be configured with PSU-managed software before use. This includes required security tools, network access, and system configurations. To complete this setup, take your new computer to the IT service desk located in 401 Ag Administration Building, Monday through Friday, from 8:00 am–12:00 pm and 1:00–5:00 pm.

You will need to drop off the computer; setup typically takes between 8 and 48 hours. Before proceeding, ensure that any personal files are backed up, as the system configuration may require administrative changes.

Once the computer is returned, you should be able to log in using your PSU credentials and access core services (email, Canvas, LionPATH, VPN, and network drives).

For research and computational work, I recommend requesting local administrator privileges on Mac computers. This is often necessary to install scientific software, manage libraries, and update dependencies. To request admin access, contact Penn State IT support (https://www.it.psu.edu/support/) and submit the required administrator access request form.

\newline

https://forms.office.com/pages/responsepage.aspx?id=RY30fNs9iUOpwcEVUm61Ltud-7yJr19Bnsvi3eO1unBUMEU5UjVUVklMNkJNTDk1NFVIWDRZSFAyMy4u&route=shorturl

\subsection{Trainings}
All students are required to complete mandatory Penn State trainings before engaging in coursework, research, or laboratory activities. Depending on your appointment type, these trainings are accessed through Workday/WorkLion (https://worklion.psu.edu/). You may be assigned several hours of online coursework, each followed by required assessments. These trainings must be completed and passed by their assigned deadlines.

Training completion can take several days. Do not rush through the material, but plan accordingly to avoid delays in your academic or research activities. Failure to complete required training may restrict access to facilities or systems. For example, the laboratory, greenhouse facilities, etc.

Laboratory access is contingent upon completion of all required safety training. After completing laboratory safety training (links will be provided by Ed), you will be eligible to obtain keys and access to lab spaces.

Penn State enforces strict compliance with training requirements. Core trainings typically include responsible conduct of research, academic integrity, and laboratory and environmental health and safety. Your academic performance and adherence to laboratory safety protocols are my highest priorities.


\subsection{Classes}
You can access your course schedule, enrollment information, and registration tools through LionPATH (https://www.lionpath.psu.edu/). All course materials, announcements, assignments, and grades are typically managed through Canvas (https://www.it.psu.edu/services/canvas/). You are expected to check Canvas and your Penn State email regularly, as instructors use these platforms for official communication.

Attendance expectations, course modality (synchronous, asynchronous, in-person, hybrid, or online), and grading policies are defined by each instructor and outlined in the course syllabus. Office hours and instructor contact information are also available on Canvas and should be consulted proactively if questions arise. You need to obtain a B (83-87\%) or above.

Penn State uses Starfish (https://sites.psu.edu/starfishinfo/) to provide academic feedback and notifications about support. You may receive alerts or evaluations at the midpoint and end of the semester; these are intended to help you track your progress and address potential concerns early.

Be aware of important academic deadlines (add/drop, late registration, withdrawal) listed on the Penn State academic calendar. It is your responsibility to monitor these dates and ensure your enrollment is accurate.


\subsection{Syllabus}
In the U.S. academic system, each course provides a document called a \textit{syllabus}. The syllabus outlines the course organization and expectations for students throughout the semester. It includes the course description, learning objectives, weekly topics, instructional format, grading and assessment criteria, required materials, key policies, and important dates.

The syllabus functions as an official contract between the instructor and the students. It defines responsibilities on both sides, establishes evaluation standards, and explains how grades are determined. Students are expected to read the syllabus carefully at the beginning of the semester and refer to it regularly, as it governs course requirements, deadlines, and academic expectations.


\section{Living in State College}

\subsection{Transportation}
State College and the University Park campus are served by the Centre Area Transportation Authority (CATA) bus system. CATA buses provide extensive coverage across campus, downtown State College, and surrounding residential areas.

Penn State graduate students can ride CATA buses with a RIDE pass. You can enroll, and it will be deducted from your payroll. Single fares are $2.50, and a daypass $5. You should verify your status and register directly with Penn State Transportation Services (Offices in the Eisenhower Deck, aka parking lot).

Up-to-date information on bus routes, fares, and registration is available at:
\begin{itemize}
  \item Penn State Transportation Services: https://transportation.psu.edu/
  \item CATA: https://catabus.com/
\end{itemize}

Download the following apps: Token Transit for payments, and myStop.

\subsection{Weather expectations}
Central Pennsylvania has a humid continental climate with precipitation distributed throughout the year. Seasonal temperature variation is pronounced.

During summer (June--August), mean temperatures in State College are approximately 22--24~$^\circ$C (72--75~$^\circ$F), with daytime highs commonly reaching 27--30~$^\circ$C (80--86~$^\circ$F). Humidity can be moderate to high.

During winter (December--February), mean temperatures are approximately $-4$ to 1~$^\circ$C (28--34~$^\circ$F). Nighttime temperatures frequently fall to $-7~^\circ$C (19~$^\circ$F) or lower. Wind chill can reduce the perceived (“real feel”) temperature to approximately $-10$ to $-17~^\circ$C (14 to 1~$^\circ$F).

Warm winter clothing—including a heavy jacket, gloves, a hat, and waterproof footwear—is strongly recommended. Snow and icy conditions are common in winter and can affect commuting and daily activities. Sign up for the University text system to get alerts.

\subsection{Grocery shopping}
State College has several grocery and convenience store options. The most common and generally affordable option is Walmart, the only full Walmart store in the immediate area, which offers groceries, household supplies, and basic pharmacy services. Additional major grocery stores include Giant and Weis, both of which are good options for routine grocery shopping.

Other grocery and specialty stores in the area include Trader Joe’s, Aldi, Wegmans, and Target. These stores vary in price, product selection, and emphasis (e.g., international foods, organic products, or prepared meals), and many students choose a combination of stores depending on their needs.

Additionally, there are International and Asian markets; unfortunately, State College lacks a Latin or Mexican store, the closest one is in Harrisburg, named Los 3 Hermanos.

Pharmacy services are available at standalone locations, such as CVS, and within larger stores, such as Walmart and Target. Most pharmacies carry over-the-counter medications and personal care products and offer prescription services.

\subsection{Manual Still under construction, but I hope it will answer a couple of questions.}

\end{document}  